%=======================================================================
%  GuardianLink -- Capstone Project Document
%  Compile:  pdflatex GuardianLink-Capstone.tex   (run twice for TOC)
%=======================================================================
\documentclass[11pt,letterpaper]{article}

\usepackage[T1]{fontenc}
\usepackage[utf8]{inputenc}
\usepackage[margin=1in]{geometry}
\usepackage{lmodern}
\usepackage{microtype}
\usepackage{booktabs}
\usepackage{tabularx}
\usepackage{longtable}
\usepackage{array}
\usepackage{enumitem}
\usepackage{fancyvrb}
\usepackage{float}
\usepackage{fancyhdr}
\usepackage{titlesec}
\usepackage{xcolor}
\usepackage[colorlinks=true,linkcolor=black,urlcolor=blue,
            citecolor=black]{hyperref}

%----------------------------- styling --------------------------------
\definecolor{glnavy}{RGB}{20,40,80}

\titleformat{\section}{\Large\bfseries\color{glnavy}}{\thesection}{1em}{}
\titleformat{\subsection}{\large\bfseries}{\thesubsection}{1em}{}
\titleformat{\subsubsection}{\normalsize\bfseries}{\thesubsubsection}{1em}{}
\titlespacing*{\section}{0pt}{18pt}{8pt}

\pagestyle{fancy}
\fancyhf{}
\lhead{\small GuardianLink --- Capstone Project Document}
\rhead{\small\thepage}
\renewcommand{\headrulewidth}{0.4pt}

\setlist[itemize]{itemsep=2pt,topsep=4pt,parsep=0pt}
\setlist[enumerate]{itemsep=2pt,topsep=4pt,parsep=0pt}

\newcolumntype{L}[1]{>{\raggedright\arraybackslash}p{#1}}
\newcommand{\wf}[1]{\begin{figure}[H]\centering
  \BVerbatimInput{#1}\end{figure}}

%--------------------------- title page -------------------------------
\begin{document}

\begin{titlepage}
\centering
\vspace*{5cm}
{\Huge\bfseries\color{glnavy} GuardianLink\par}
\vspace{0.6cm}
{\LARGE Capstone Project Document\par}
\vspace{2.5cm}
{\large
\begin{tabular}{rl}
\textbf{Student:}        & [Your Name]\\[6pt]
\textbf{Submission date:}& [Date]\\[6pt]
\textbf{Repository:}     & \url{https://github.com/your-username/guardianlink}\\[6pt]
\textbf{Walk-through:}   & \url{https://jmp.sh/your-link}\\
\end{tabular}\par}
\vfill
{\small A full-stack web platform connecting volunteer cybersecurity\\
specialists with non-profit organizations, at no cost to the non-profit.\par}
\vspace{1cm}
\end{titlepage}

\tableofcontents
\newpage

%=======================================================================
\section{Overview and Summary}
%=======================================================================

\subsection{Synopsis}

GuardianLink is a full-stack web platform that connects volunteer
cybersecurity specialists with non-profit organizations and NGOs that need
security assistance but cannot afford commercial consulting rates. The
platform allows both groups to register, maintain a profile, browse the
opposite user type, and initiate contact --- all at no cost to the
non-profit.

The application is a single Node.js project: an Express REST API backed by
MongoDB, serving a React single-page frontend. Authentication is JWT-based
with role-based access control enforcing that NGOs and volunteers can only
discover each other, never their own kind.

\subsection{Purpose}

Non-profits hold sensitive donor, client, and beneficiary data but rarely
have the budget for a security audit, policy review, or incident-response
plan. Meanwhile, many working security professionals are willing to donate a
few hours a week but have no reliable channel for finding organizations that
need them. GuardianLink removes the discovery problem on both sides.

\subsection{Goals}

\begin{center}
\begin{tabularx}{\textwidth}{@{}c X l@{}}
\toprule
\textbf{\#} & \textbf{Goal} & \textbf{Status}\\
\midrule
1  & Public landing page explaining who we are, what we do, why, and our partners & Complete\\
2  & Separate registration flows capturing NGO and volunteer fields & Complete\\
3  & Three enforced user roles: NGO, Volunteer, Admin & Complete\\
4  & Login required to view any profile & Complete\\
5  & Cross-role visibility only (NGO / Volunteer) & Complete\\
6  & NGO-to-volunteer communication channel & Complete (mailto)\\
7  & Forgot-password flow & Complete\\
8  & Admin can delete users, assign roles, reset passwords, create profiles & Complete\\
9  & At least one admin exists at all times & Complete\\
10 & Users can delete their own account, and only their own & Complete\\
11 & Responsive desktop and mobile layout & Complete\\
\bottomrule
\end{tabularx}
\end{center}

\subsection{Requirements Coverage}

\subsubsection*{Information collected from NGOs / Non-Profits}
\begin{itemize}
  \item Point-of-contact email
  \item Organization name
  \item Areas of concern
  \item Password, plus an optional description field
\end{itemize}

\subsubsection*{Information collected from Cybersecurity Volunteers}
\begin{itemize}
  \item Name
  \item Email
  \item Hours available per week
  \item Criminal background check status
  \item LinkedIn profile URL (serves as the resume)
  \item Areas of expertise, plus an optional bio
\end{itemize}

\subsubsection*{User types}
\begin{enumerate}
  \item \textbf{NGO / Small Company} --- registers via the Organization
        Registration form; browses the volunteer directory; contacts volunteers.
  \item \textbf{Volunteer (Cybersecurity Specialist)} --- registers via the
        Volunteer Registration form; browses the NGO directory; responds to requests.
  \item \textbf{Admin} --- manages all users and oversees platform operations.
        The system refuses any operation that would leave zero admins.
\end{enumerate}

%=======================================================================
\section{Wireframe Diagrams and Planning Documents}
%=======================================================================

\subsection{Site Map}

\begin{figure}[H]
\centering
\begin{BVerbatim}[fontsize=\footnotesize]
                        +--------------+
                        |   Landing    |   (public)
                        |    / Home    |
                        +------+-------+
             +-----------------+-----------------+
             |                 |                 |
             v                 v                 v
      +------------+    +------------+    +------------+
      |   About    |    |  Register  |    |   Login    |
      |  (public)  |    |  (public)  |    |  (public)  |
      +------------+    +-----+------+    +-----+------+
                              |                 |
                   +----------+------+          |   +-----------------+
                   |                 |          |   | Forgot Password |
                   v                 v          |   |     (modal)     |
            +-------------+  +--------------+   |   +-----------------+
            |     NGO     |  |  Volunteer   |   |
            |  Register   |  |   Register   |   |
            +------+------+  +------+-------+   |
                   +--------+-------+           |
                            |                   |
                            v                   |
          ============ AUTH WALL ===============+=============
                            |                   |
                            v                   v
                  +------------------------------+
                  |          Directory           |
                  |     (opposite role only)     |
                  +---------------+--------------+
                                  |
                                  v
                  +------------------------------+
                  |        Profile Detail        |
                  |    [ Contact via Email ]     |
                  +------------------------------+

             +----------------+----------------+
             v                v                v
      +------------+   +------------+   +----------------+
      |  Account   |   | Edit Prof. |   |  Admin Panel   |
      | (password, |   |            |   |  (admin only)  |
      |  delete)   |   |            |   +----------------+
      +------------+   +------------+
\end{BVerbatim}
\caption{Application site map. Everything below the auth wall requires a valid JWT.}
\end{figure}

\subsection{Wireframe --- Landing Page}

\begin{figure}[H]
\centering
\begin{BVerbatim}[fontsize=\footnotesize]
+--------------------------------------------------------------+
|  GuardianLink        About   How It Works   [Log in] [Join]  |
+--------------------------------------------------------------+
|                                                              |
|         Free cybersecurity help for non-profits.             |
|         Donated by professionals who care.                   |
|                                                              |
|         [ I'm a Non-Profit ]   [ I'm a Volunteer ]           |
|                                                              |
+--------------------------------------------------------------+
|  WHO WE ARE          WHAT WE DO           WHY WE DO IT       |
|  +------------+      +------------+       +------------+     |
|  | A network  |      | Match NGOs |       | Non-profits|     |
|  | of security|      | with vetted|       | hold       |     |
|  | specialists|      | specialists|       | sensitive  |     |
|  | who donate |      | at no cost |       | data with  |     |
|  | their time |      |            |       | no budget  |     |
|  +------------+      +------------+       +------------+     |
+--------------------------------------------------------------+
|  OUR PARTNERS                                                |
|  [ Logo ]  [ Logo ]  [ Logo ]  [ Logo ]                      |
+--------------------------------------------------------------+
|  (c) GuardianLink          Privacy  |  Contact               |
+--------------------------------------------------------------+
\end{BVerbatim}
\caption{Landing page covering all four required About subsections.}
\end{figure}

\subsection{Wireframe --- Registration (role-switched form)}

\begin{figure}[H]
\centering
\begin{BVerbatim}[fontsize=\footnotesize]
+--------------------------------------------+
|            Create your account             |
|                                            |
|   ( o ) Non-Profit      (   ) Volunteer    |
|   --------------------------------------   |
|                                            |
|   Organization Name  [__________________]  |
|   Contact Email      [__________________]  |
|   Password           [__________________]  |
|   Areas of Concern   [ ] Network Security  |
|                      [ ] Phishing / Email  |
|                      [ ] Policy & Training |
|                      [ ] Incident Response |
|   Short Description  [__________________]  |
|                      [__________________]  |
|                                            |
|            [   Create Account   ]          |
|       Already have an account? Log in      |
+--------------------------------------------+

    When "Volunteer" is selected, the lower
    half swaps to:

|   Full Name          [__________________]  |
|   Email              [__________________]  |
|   Password           [__________________]  |
|   Hours / week       [____]   (1-60)       |
|   Areas of Expertise [ ] Pen Testing ...   |
|   LinkedIn URL       [__________________]  |
|   [X] I consent to a criminal background   |
|       check                                |
|   Short Bio          [__________________]  |
\end{BVerbatim}
\caption{Single registration component; the role radio swaps the field set.}
\end{figure}

\subsection{Wireframe --- Directory (post-login)}

\begin{figure}[H]
\centering
\begin{BVerbatim}[fontsize=\footnotesize]
+--------------------------------------------------------------+
| GuardianLink    Directory   Account   [Admin]      Log out   |
+--------------------------------------------------------------+
|  Volunteers                        Search [____________]     |
|  Filter: [ All areas  v ]   [ ] Background-checked only      |
+--------------------------------------------------------------+
|  +--------------------+  +--------------------+              |
|  | Amara Okafor       |  | Devin Reyes        |              |
|  | Pen Testing, IR    |  | Policy, Training   |              |
|  | 10 hrs/week        |  | 5 hrs/week         |              |
|  | [X] Background chk |  | [ ] Not checked    |              |
|  | [View] [Contact]   |  | [View] [Contact]   |              |
|  +--------------------+  +--------------------+              |
|                                                              |
|  +--------------------+  +--------------------+              |
|  | ...                |  | ...                |              |
|  +--------------------+  +--------------------+              |
+--------------------------------------------------------------+

   An NGO account never sees NGO cards here; a volunteer
   account sees only NGO cards. Enforced server-side.
\end{BVerbatim}
\caption{Role-filtered directory. Visibility is derived from the JWT.}
\end{figure}

\subsection{Wireframe --- Admin Panel}

\begin{figure}[H]
\centering
\begin{BVerbatim}[fontsize=\footnotesize]
+--------------------------------------------------------------+
|  Admin Panel                            [ + Create User ]    |
+--------------------------------------------------------------+
|  Name          Email              Role       Actions         |
| ------------------------------------------------------------ |
|  Site Admin    admin@guardi...    Admin v   [Reset] [Del]    |
|  Tom Alvarez   tom@harvest...     NGO   v   [Reset] [Del]    |
|  Amara Okafor  amara.sec@...      Vol.  v   [Reset] [Del]    |
+--------------------------------------------------------------+
|  PASSWORD RESET REQUESTS                                     |
|  tom@harvesttable.org    2 hours ago    [Resolve] [Dismiss]  |
+--------------------------------------------------------------+

   Attempting to delete or demote the last remaining admin
   returns 409 Conflict and the row is left unchanged.
\end{BVerbatim}
\caption{Admin panel with the last-admin guard.}
\end{figure}

\subsection{Project Plan --- Build Order}

Components were built bottom-up so each layer could be tested before the
next depended on it.

\subsubsection*{Component 1 --- Data Model and Database}
\textit{Subtasks}
\begin{itemize}
  \item Install and connect MongoDB (local service, later MongoDB Atlas)
  \item Define the \texttt{User} schema with role-discriminating fields
  \item Define the \texttt{ResetRequest} schema
  \item Add a \texttt{pre("save")} hook hashing passwords with bcrypt at 12 rounds
  \item Add a \texttt{comparePassword} instance method
  \item Add a \texttt{toSafeJSON} method stripping the hash and mapping
        \texttt{\_id} to \texttt{id}
  \item Write \texttt{seed.js} creating one admin, one NGO, and one volunteer
\end{itemize}
\textit{Acceptance criteria}
\begin{itemize}
  \item \texttt{node seed.js} completes without error and creates exactly three users
  \item Every stored password matches the bcrypt pattern below and is 60 characters long
  \item No API response anywhere contains a \texttt{password} field
  \item Re-running the seed is idempotent --- it does not produce duplicates
\end{itemize}
\begin{Verbatim}[fontsize=\small,xleftmargin=1.5em]
/^\$2[aby]\$12\$/
\end{Verbatim}

\subsubsection*{Component 2 --- Authentication}
\textit{Subtasks}
\begin{itemize}
  \item \texttt{POST /api/auth/register} for both roles
  \item \texttt{POST /api/auth/login} issuing a signed JWT
  \item \texttt{GET /api/auth/me} returning the current user
  \item \texttt{requireAuth} middleware verifying the bearer token
  \item \texttt{passwordChangedAt} comparison against the token \texttt{iat} claim
\end{itemize}
\textit{Acceptance criteria}
\begin{itemize}
  \item Correct credentials return a token and a safe user object
  \item Wrong password and unknown email return an identical 401 message, so the
        endpoint cannot be used to enumerate registered emails
  \item A request with no token, a malformed token, or an expired token returns 401
  \item A token minted before a password change is rejected on the next request
\end{itemize}

\subsubsection*{Component 3 --- Role-Based Access Control}
\textit{Subtasks}
\begin{itemize}
  \item \texttt{requireRole(...roles)} middleware
  \item \texttt{GET /api/directory} filtering by the opposite role
  \item Derive visibility from the JWT, never from a client-supplied query parameter
\end{itemize}
\textit{Acceptance criteria}
\begin{itemize}
  \item An NGO token returns only volunteers; a volunteer token only NGOs;
        an admin token everyone
  \item Tampering with a request to ask for same-role profiles still returns
        the correct filtered set
  \item Admin-only routes return 403 for NGO and volunteer tokens
\end{itemize}

\subsubsection*{Component 4 --- Public Pages and Registration UI}
\textit{Subtasks}
\begin{itemize}
  \item Landing page with Who / What / Why / Partners sections
  \item Role-switching registration form
  \item Client-side validation before submit
  \item Responsive layout down to 375\,px
\end{itemize}
\textit{Acceptance criteria}
\begin{itemize}
  \item All four About subsections are present and readable
  \item Switching the role radio swaps the correct field set
  \item Every required NGO and volunteer field from the brief is captured and persisted
  \item No horizontal scrolling at 375\,px, 768\,px, or 1440\,px
\end{itemize}

\subsubsection*{Component 5 --- Directory and Profile Views}
\textit{Subtasks}
\begin{itemize}
  \item Directory grid with search and filter
  \item Profile detail view
  \item \texttt{mailto:} contact button with prefilled recipient, subject, and body
\end{itemize}
\textit{Acceptance criteria}
\begin{itemize}
  \item Visiting the directory while logged out redirects to login
  \item The contact button opens the default mail client with all three fields populated
  \item Search filters by name, organization, and area of concern
\end{itemize}

\subsubsection*{Component 6 --- Account Self-Service}
\textit{Subtasks}
\begin{itemize}
  \item Profile editing
  \item Self-service password change requiring the current password
  \item Self-deletion
\end{itemize}
\textit{Acceptance criteria}
\begin{itemize}
  \item Changing a password without the correct current password fails with 401
  \item After a successful change, the old password no longer works and the
        old token is rejected
  \item A user can delete their own account and no other
\end{itemize}

\subsubsection*{Component 7 --- Forgot Password}
\textit{Subtasks}
\begin{itemize}
  \item Public request form creating a \texttt{ResetRequest}
  \item Admin queue listing outstanding requests
  \item Admin force-set that bypasses the current-password requirement
\end{itemize}
\textit{Acceptance criteria}
\begin{itemize}
  \item Submitting a request for an unknown email returns the same confirmation
        as a known one, revealing nothing
  \item The request appears in the admin queue
  \item After an admin reset, the user can log in with the new password
\end{itemize}

\subsubsection*{Component 8 --- Admin Panel}
\textit{Subtasks}
\begin{itemize}
  \item User table with role dropdown and delete action
  \item Create-user form
  \item Password reset modal with \texttt{isSelf} detection
  \item Last-admin guards on delete and demote
\end{itemize}
\textit{Acceptance criteria}
\begin{itemize}
  \item Deleting or demoting the only admin returns 409 and changes nothing
  \item Creating a user does not invalidate the admin's own session
  \item An admin resetting their own password is prompted for the current one;
        resetting another user's password is not
\end{itemize}

\subsubsection*{Component 9 --- Hardening and Deployment}
\textit{Subtasks}
\begin{itemize}
  \item \texttt{helmet}, \texttt{express-rate-limit}, \texttt{express-mongo-sanitize}
  \item Centralised error handler
  \item Migrate to MongoDB Atlas
  \item Deploy to Render binding \texttt{0.0.0.0}
\end{itemize}
\textit{Acceptance criteria}
\begin{itemize}
  \item Repeated failed logins are throttled
  \item A NoSQL injection payload in the login body fails to authenticate
  \item No stack traces are returned to clients in production
  \item The deployed URL serves the app and all routes behave as they do locally
\end{itemize}

\subsection{Acceptance Criteria --- Whole Project}

The project is complete when all of the following hold on the deployed instance:

\begin{enumerate}
  \item An unauthenticated visitor can read the About content but reach no
        profile data by any route.
  \item An NGO and a volunteer can each register with every field required by
        the brief, and those values are correctly persisted and redisplayed.
  \item Each role sees only the opposite role in the directory, enforced on the server.
  \item An NGO can open a prefilled email to a volunteer in one click.
  \item A user who has forgotten their password can regain access through the
        admin-mediated flow.
  \item An admin can create, delete, promote, demote, and reset the password of any user.
  \item The system refuses every operation that would leave the platform with no admin.
  \item A regular user can delete their own account and cannot affect any other account.
  \item Passwords are never stored in plaintext, never returned by the API, and never logged.
  \item The UI is usable on both a phone and a desktop.
  \item The full source is on GitHub with a meaningful commit history.
  \item A walk-through video demonstrates every Must Have.
\end{enumerate}

%=======================================================================
\section{Technology Stack and Rationale}
%=======================================================================

\begin{center}
\begin{tabularx}{\textwidth}{@{}L{2.2cm} L{3.4cm} X@{}}
\toprule
\textbf{Layer} & \textbf{Technology} & \textbf{Why}\\
\midrule
Runtime & Node.js &
One language across the whole stack; no context switching between server and browser\\
\addlinespace
Framework & Express &
Minimal and unopinionated; middleware chaining maps directly onto the
auth $\rightarrow$ role $\rightarrow$ handler pipeline the brief requires\\
\addlinespace
Database & MongoDB + Mongoose 8.5.1 &
NGO and volunteer records hold genuinely different fields. A document store
absorbs that without nullable columns or a join table. Mongoose adds the
schema validation and hooks a raw driver lacks\\
\addlinespace
Frontend & React 18 via CDN, Babel 7.26.4 &
Component model and state handling without a build step. The entire frontend
is one \texttt{index.html}, so there is no bundler to misconfigure\\
\addlinespace
Auth & \texttt{jsonwebtoken} &
Stateless sessions. The token carries the role, so no session store is needed
and horizontal scaling stays trivial\\
\addlinespace
Hashing & \texttt{bcryptjs}, 12 rounds &
Pure JavaScript, so no native build toolchain on Windows. Twelve rounds is the
current practical balance of cost and security\\
\addlinespace
Security & \texttt{helmet}, \texttt{express-rate-limit}, \texttt{express-mongo-sanitize} &
Sensible headers, brute-force throttling, and operator-injection stripping ---
appropriate for a platform whose subject matter is security\\
\addlinespace
Config & \texttt{dotenv} &
Keeps \texttt{MONGO\_URI} and \texttt{JWT\_SECRET} out of version control\\
\addlinespace
Dev tooling & \texttt{nodemon}, WebStorm &
Auto-restart on save; the IDE run console surfaces stack traces with
clickable file links\\
\addlinespace
Hosting & Render + MongoDB Atlas M0 &
Both have free tiers adequate for a capstone demo and connect over a
standard URI\\
\bottomrule
\end{tabularx}
\end{center}

\subsection{Why a single Node project}

An earlier iteration ran the frontend as a separate static project against a
mock API. Collapsing it into one Express app that serves
\texttt{public/index.html} eliminated CORS configuration, gave one command to
start everything, and made deployment a single service rather than two.

\subsection{Why JWT rather than server sessions}

The access-control rule is role-driven, so the role needs to be available on
every request. Embedding the user id in a signed token and loading the user in
\texttt{requireAuth} keeps that lookup in one place. The known weakness ---
tokens staying valid until expiry --- is addressed by the
\texttt{passwordChangedAt} check described below.

%=======================================================================
\section{Testing Methods}
%=======================================================================

\subsection{Black Box Testing}

The full RBAC matrix was exercised through the UI with no reference to the
implementation.

\begin{center}
\begin{tabular}{@{}llll@{}}
\toprule
\textbf{Logged in as} & \textbf{Sees in directory} & \textbf{Admin panel} & \textbf{Can delete others}\\
\midrule
NGO       & Volunteers only & Hidden  & No\\
Volunteer & NGOs only       & Hidden  & No\\
Admin     & All users       & Visible & Yes\\
\bottomrule
\end{tabular}
\end{center}

Additional black-box cases: registering with an already-used email, logging in
with a wrong password, submitting a forgot-password request for an address
that does not exist, and submitting the volunteer form with an empty hours field.

\subsection{White Box Testing}

Tested against knowledge of the internals, targeting specific branches:

\begin{itemize}
  \item \textbf{Hash-state inspection.} A \texttt{check.js} script connects
        directly to MongoDB, prints the stored hash, asserts it matches the
        bcrypt prefix pattern, and runs \texttt{bcrypt.compare} against the
        expected plaintext. This distinguishes a bad login route from a bad
        hash, which look identical from outside.
  \item \textbf{Token-vs-timestamp branch.} Captured a token, changed the
        password, replayed the token, and confirmed the
        \texttt{payload.iat < passwordChangedAt} branch rejected it.
  \item \textbf{Pre-save hook guard.} Verified that
        \texttt{isModified("password")} prevents re-hashing when only a role or
        active flag is saved.
  \item \textbf{Last-admin guard.} Reduced the database to a single admin and
        confirmed both the delete and the demote paths reach the 409 branch.
\end{itemize}

\subsection{Unit Testing}

Isolated checks on individual functions:

\begin{itemize}
  \item \texttt{comparePassword} returns \texttt{true} for the correct password,
        \texttt{false} otherwise
  \item \texttt{toSafeJSON} omits \texttt{password} and exposes \texttt{id}
        rather than \texttt{\_id}
  \item The registration normaliser converts a comma-separated string to a clean
        array, coerces checkbox values of \texttt{"on"}, \texttt{"true"}, and
        \texttt{true} to boolean \texttt{true}, and omits \texttt{hours}
        entirely when the input is blank
  \item \texttt{requireRole} returns 403 for a role not in its allowed list
\end{itemize}

\subsection{Integration Testing}

Endpoints were exercised with \texttt{curl.exe} against a running server,
chaining the token from \texttt{/api/auth/login} into an
\texttt{Authorization: Bearer} header for subsequent calls. This caught faults
the browser masked --- most usefully, it distinguished ``the server returned
401'' from ``nothing is listening on the port,'' which present identically in
a browser.

\subsection{Regression Testing}

After every significant change, a fixed checklist was re-run: seed, log in as
each of the three roles, load the directory, open a profile, send a contact
email, change a password, and perform one admin action. This caught the case
where a change-password fix broke login for accounts that had never changed
their password.

\subsection{Security Testing}

\begin{itemize}
  \item Submitted the payload below to the login endpoint to confirm
        \texttt{express-mongo-sanitize} strips the operators and authentication fails
  \item Sent \texttt{"role": "admin"} in a registration payload to confirm the
        explicit field mapping ignores it
  \item Repeated failed logins to confirm the rate limiter engages
  \item Inspected every API response in DevTools to confirm no password hash is
        ever transmitted
\end{itemize}

\begin{Verbatim}[fontsize=\small,xleftmargin=1.5em]
{"email": {"$gt": ""}, "password": {"$gt": ""}}
\end{Verbatim}

%=======================================================================
\section{Challenges Faced}
%=======================================================================

\subsection{Code Pasted Into the Wrong File}

The single most costly problem. A login query was pasted into
\texttt{models/User.js}, replacing the schema definition. Node reported it as
a misleading ES-module warning about setting \texttt{"type": "module"}. The
real fault was a top-level \texttt{await} in a CommonJS file, which Node
guesses must be an ESM problem. Because the model no longer exported anything
usable, the seed script failed, the login endpoint failed, and the diagnostic
script failed --- three unrelated-looking symptoms from one cause. The lesson
was to run \texttt{node -{}-check} across every file before interpreting a
confusing error message.

\subsection{Double-Hashed Passwords}

Login rejected credentials that were definitely correct. The stored value
looked like a valid bcrypt hash, so nothing seemed wrong on inspection. The
seed script was hashing the password manually \textit{and} the
\texttt{pre("save")} hook was hashing it again. The fix was to designate the
hook as the only place hashing ever happens and pass plaintext everywhere
else. A related latent bug was found at the same time: without an
\texttt{isModified("password")} guard, every unrelated \texttt{save()} ---
such as changing a role --- silently re-hashed the stored hash and destroyed
the account's password.

\subsection{Missing Middleware Import}

\texttt{ReferenceError: requireAuth is not defined} --- the middleware was
defined and exported, but a route file never imported it. The error class
mattered: a \texttt{ReferenceError} means the identifier does not exist in
scope, whereas a bad export shape produces \texttt{Route.get() requires a
callback function but got [object Undefined]}. Distinguishing the two narrowed
the search from ``somewhere in auth'' to ``the top of one file.''

\subsection{Server Failing to Bind}

The browser reported \texttt{ERR\_CONNECTION\_REFUSED} while the terminal
appeared to have started the server. Two causes compounded:
\texttt{app.listen} sat behind an unhandled \texttt{mongoose.connect}
rejection, so a database failure killed the process with no message; and the
server was being started in the same terminal window later used for
\texttt{curl}, where Ctrl+C had stopped it. Restructuring startup so
\texttt{app.listen} runs inside \texttt{.then()} with an explicit
\texttt{.catch}, plus an \texttt{unhandledRejection} handler, turned silent
exits into readable errors.

\subsection{Volunteer Fields Not Persisting}

After registration, \texttt{areas} saved as an empty array,
\texttt{backgroundCheck} was always \texttt{false}, and \texttt{hours} was
absent from the document entirely. The pattern was diagnostic: every field
with a schema default received its default, and the one field without a
default vanished --- the signature of values never reaching Mongoose at all.
The register handler destructured a fixed list of keys and dropped the rest. A
second defect compounded it: the form sent \texttt{background} while the schema
declared \texttt{backgroundCheck}, and Mongoose's strict mode discards unknown
keys silently rather than erroring.

\subsection{Method Name Drift}

\texttt{TypeError: user.checkPassword is not a function} --- the model defined
\texttt{comparePassword} and a route called \texttt{checkPassword}. Trivial
once seen, but it highlighted a structural risk: the same concept had acquired
two names across files, exactly as \texttt{orgName} and \texttt{org} had on the
frontend boundary.

\subsection{Session Invalidation Side Effects}

Adding \texttt{passwordChangedAt} correctly invalidated old tokens, which then
produced confusing 401 errors immediately after a successful password change
--- the frontend was still holding the now-dead token. The fix was to clear
stored credentials and return the user to sign-in on success.

\subsection{Migration From Mock to Real Backend}

The frontend was originally built against a \texttt{localStorage} mock whose
method names did not match the eventual REST routes. Rather than rewrite every
call site, the mock's interface was preserved as a thin \texttt{fetch}-based
client that maps old names to new endpoints and normalises field names in both
directions. Dead mock code was then removed. Leftover hardcoded demo
credentials in the UI caused a separate false alarm when they no longer matched
what \texttt{seed.js} actually created.

%=======================================================================
\section{Lessons Learned}
%=======================================================================

\begin{description}[style=unboxed,leftmargin=0pt,itemsep=6pt]
  \item[The error message often names the wrong problem.] Node's ESM warning
    pointed at module configuration when the actual fault was misplaced code.
    Reading the line number and the quoted source line proved far more reliable
    than reading the message text.
  \item[Check the data before debugging the code.] Writing a short script to
    print the stored hash and run \texttt{bcrypt.compare} directly against it
    answered in seconds what would otherwise have taken an hour of reading
    route logic. Inspecting state at the boundary is almost always faster than
    reasoning about the layers above it.
  \item[The shape of a failure narrows the cause.] Fields with defaults
    receiving their defaults while a field without one disappeared entirely was
    conclusive evidence the data never arrived --- no logging needed to rule
    out a validation failure.
  \item[Designate exactly one place for each responsibility.] Hashing in both
    the seed and the model hook caused the double-hash bug. Making the
    \texttt{pre("save")} hook the sole owner removed an entire class of error.
  \item[Guard hooks that mutate.] A \texttt{pre("save")} hook without an
    \texttt{isModified} check is a latent data-destruction bug that will not
    surface until an unrelated feature calls \texttt{save()}.
  \item[Name one concept once.] \texttt{checkPassword} versus
    \texttt{comparePassword}, \texttt{background} versus
    \texttt{backgroundCheck}, \texttt{org} versus \texttt{orgName} --- three
    bugs with the same root cause. Where a boundary genuinely needs two names,
    the translation belongs in one explicit place.
  \item[Never pass \texttt{req.body} straight to a model constructor.] Mapping
    fields explicitly is slightly more code and closes a privilege-escalation
    hole where a crafted payload could set \texttt{role: "admin"}.
  \item[Test the layer below before blaming the layer above.] A browser cannot
    distinguish ``server returned an error'' from ``no server is running.'' A
    \texttt{curl} exit code can.
  \item[Commit more often.] Several of these faults would have been a one-line
    \texttt{git diff} had commits been smaller.
\end{description}

%=======================================================================
\section{Future Recommendations}
%=======================================================================

\subsection{Security}
\begin{itemize}
  \item Replace admin-mediated password reset with emailed single-use tokens,
        hashed at rest with a 15-minute expiry
  \item Add refresh tokens with a short-lived access token, moving them to
        \texttt{httpOnly} cookies to remove the XSS exposure of
        browser-stored JWTs
  \item Offer TOTP two-factor authentication, particularly for admin accounts
  \item Require email verification at registration
  \item Add an append-only audit log of admin actions --- deletions, role
        changes, password resets
\end{itemize}

\subsection{Functionality}
\begin{itemize}
  \item In-app messaging to replace \texttt{mailto:}, giving a conversation
        history inside the platform
  \item A structured engagement workflow: request $\rightarrow$ accept
        $\rightarrow$ in progress $\rightarrow$ complete, with hours logged
        against each
  \item Volunteer availability calendar and matching suggestions based on
        declared areas
  \item Verified background-check status reviewed by an admin, rather than
        self-declared
  \item Ratings or references after a completed engagement
  \item Admin dashboard with registration, match, and volunteer-hour metrics
\end{itemize}

\subsection{Engineering}
\begin{itemize}
  \item Introduce a build step with Vite, removing the in-browser Babel
        transform and its runtime cost
  \item Automated test suite with Jest and Supertest covering the RBAC matrix
        and auth edge cases, run in CI on every push
  \item Split \texttt{index.html} into separate component modules as the
        frontend grows
  \item Centralise validation with a schema validator such as Zod, shared
        between client and server, which would have prevented the
        field-mapping bug outright
  \item Structured request logging with Pino
  \item Database indexes on \texttt{role} and \texttt{areas} to keep directory
        queries fast at scale
  \item Pagination on the directory endpoint
\end{itemize}

\subsection{Accessibility and UX}
\begin{itemize}
  \item Full WCAG 2.1 AA audit: keyboard navigation, focus management in
        modals, colour contrast
  \item Screen-reader labelling on all form controls and ARIA live regions
        for toasts
  \item Inline per-field validation messages rather than a single summary error
\end{itemize}

%=======================================================================
\section{Appendix}
%=======================================================================

\subsection{Repository Structure}

\begin{figure}[H]
\centering
\begin{BVerbatim}[fontsize=\small]
GuardianLink Project/
|-- server.js        Express app, Mongo connection, static serving
|-- seed.js          Creates demo admin, NGO, and volunteer
|-- check.js         Diagnostic: inspects stored hashes
|-- package.json
|-- .env             PORT, MONGO_URI, JWT_SECRET, JWT_EXPIRES
|-- .gitignore       Excludes .env and node_modules
|-- models/
|   |-- User.js      Schema, hash hook, comparePassword, toSafeJSON
|   +-- ResetRequest.js
|-- middleware/
|   +-- auth.js      requireAuth, requireRole
|-- routes/
|   |-- auth.js      register, login, me, change-password
|   |-- users.js     self-service profile, password, delete
|   |-- directory.js role-filtered listing
|   +-- admin.js     user management, reset queue
+-- public/
    +-- index.html   React SPA
\end{BVerbatim}
\end{figure}

\subsection{API Reference}

\begin{center}
\begin{longtable}{@{}l L{6.2cm} l X@{}}
\toprule
\textbf{Method} & \textbf{Route} & \textbf{Auth} & \textbf{Purpose}\\
\midrule
\endfirsthead
\toprule
\textbf{Method} & \textbf{Route} & \textbf{Auth} & \textbf{Purpose}\\
\midrule
\endhead
POST   & \texttt{/api/auth/register}            & Public & Create an NGO or volunteer account\\
POST   & \texttt{/api/auth/login}               & Public & Authenticate, return JWT\\
GET    & \texttt{/api/auth/me}                  & Any    & Current user\\
POST   & \texttt{/api/auth/change-password}     & Any    & Self-change, requires current password\\
POST   & \texttt{/api/auth/forgot-password}     & Public & Create a reset request\\
GET    & \texttt{/api/directory}                & Any    & Opposite-role profiles\\
GET    & \texttt{/api/users/:id}                & Any    & Single profile, role-filtered\\
PATCH  & \texttt{/api/users/me}                 & Any    & Update own profile\\
DELETE & \texttt{/api/users/me}                 & Any    & Delete own account\\
GET    & \texttt{/api/admin/users}              & Admin  & All users\\
POST   & \texttt{/api/admin/users}              & Admin  & Create a user\\
PATCH  & \texttt{/api/admin/users/:id/role}     & Admin  & Change role (last-admin guarded)\\
PATCH  & \texttt{/api/admin/users/:id/password} & Admin  & Force-set a password\\
DELETE & \texttt{/api/admin/users/:id}          & Admin  & Delete a user (last-admin guarded)\\
GET    & \texttt{/api/admin/reset-requests}     & Admin  & Outstanding reset queue\\
\bottomrule
\end{longtable}
\end{center}

\subsection{Environment Variables}

\begin{Verbatim}[fontsize=\small,xleftmargin=1.5em]
PORT=3000
MONGO_URI=mongodb://127.0.0.1:27017/guardianlink
JWT_SECRET=[long random string]
JWT_EXPIRES=2h
\end{Verbatim}

\subsection{Local Setup}

\begin{Verbatim}[fontsize=\small,xleftmargin=1.5em]
git clone [repository-url]
cd "GuardianLink Project"
npm install
# create .env using the template above
node seed.js
npm run dev
# open http://localhost:3000
\end{Verbatim}

\subsection{Demo Accounts}

\begin{center}
\begin{tabular}{@{}lll@{}}
\toprule
\textbf{Role} & \textbf{Email} & \textbf{Password}\\
\midrule
Admin     & admin@guardianlink.org  & [as set in seed.js]\\
NGO       & tom@harvesttable.org    & [as set in seed.js]\\
Volunteer & amara.sec@example.com   & [as set in seed.js]\\
\bottomrule
\end{tabular}
\end{center}

\end{document}